% ============================================================
% 最小二乗法（手引書の前付け，レポートの書き方の次に入れる）
% 参考にした資料:
%   基礎ゼミ実験配布テキスト.pdf「付録A 最小二乗法」
%   レポートの書き方.pdf「2.2 最小二乗法」（原点を通る直線の場合）
% 図・表・式のラベルは lsq- で始める
% ============================================================

\introtitle{最小二乗法}

実験では，測定した2つの量の関係を式で表したい場合がよくあります。しかし，測定値には誤差があり，測定点はちょうど1本の直線や曲線の上にはのらず，少しばらつきます。このようなとき，測定値に最もよく当てはまる直線や曲線の係数（直線なら傾きと切片）を求める方法が\important{最小二乗法}です。
ここでは，抵抗器に加える電圧と流れる電流の関係のように，2つの量の関係が\important{直線}で近似できる場合の最小二乗法について説明します。

% ------------------------------------------------------------
\section{考え方}

$n$ 個の測定値 $(x_1, y_1), (x_2, y_2), \dots, (x_n, y_n)$ に，直線
\begin{equation}
  y = ax + b
  \label{eq:lsq-line}
\end{equation}
を当てはめることを考えます。各測定点について，測定値 $y_i$ と直線上の値 $ax_i+b$ との差（\cref{fig:lsq-residual}の縦の線）を\important{残差}といいます。最小二乗法では，残差を2乗して足し合わせた
\begin{equation}
  S = \sum_{i=1}^{n} \{\,y_i - (ax_i + b)\,\}^2
  \label{eq:lsq-s}
\end{equation}
が最も小さくなるように，傾き $a$ と切片 $b$ を決めます。残差を2乗するのは，直線の上側にある点（残差が正）と下側にある点（残差が負）のずれが打ち消し合わないようにするためです。

\begin{figure}[htbp]
  \centering
  \begin{tikzpicture}
    \begin{axis}[
        width=0.6\linewidth, height=0.4\linewidth,
        xmin=0, xmax=6, ymin=0, ymax=6,
        xtick={0,1,...,6}, ytick={0,1,...,6},
        xlabel={$x$}, ylabel={$y$},
        tick align=inside, axis on top,
      ]
      % 残差（測定点から直線までの縦の線）
      \addplot[EGWarn, very thick, forget plot] coordinates {(1,1.5) (1,1.380)};
      \addplot[EGWarn, very thick, forget plot] coordinates {(2,1.9) (2,2.310)};
      \addplot[EGWarn, very thick, forget plot] coordinates {(3,3.4) (3,3.240)};
      \addplot[EGWarn, very thick, forget plot] coordinates {(4,3.6) (4,4.170)};
      \addplot[EGWarn, very thick, forget plot] coordinates {(5,5.3) (5,5.100)};
      \addplot[domain=0:6, samples=2, black]{0.93*x+0.45};
      \addplot[only marks, mark=*, mark size=2pt]
        coordinates {(1,1.5) (2,1.9) (3,3.4) (4,3.6) (5,5.3)};
      \node[anchor=west, font=\small, text=EGWarn] at (axis cs:3.08,3.0) {残差};
      \node[anchor=west, font=\small] at (axis cs:0.25,5.4) {直線 $y=ax+b$};
    \end{axis}
  \end{tikzpicture}
  \caption{測定点と最小二乗法で求めた直線（縦の線が残差）}
  \label{fig:lsq-residual}
\end{figure}

% ------------------------------------------------------------
\section{傾きと切片の求め方}

\cref{eq:lsq-s}の $S$ が最小になる条件（$S$ を $a$，$b$ でそれぞれ微分して0とおく）から，傾き $a$ と切片 $b$ は次の式で求められます。
\begin{equation}
  a = \frac{n\sum x_i y_i - \sum x_i \sum y_i}{n\sum x_i^2 - \left(\sum x_i\right)^2},\qquad
  b = \frac{\sum x_i^2 \sum y_i - \sum x_i y_i \sum x_i}{n\sum x_i^2 - \left(\sum x_i\right)^2}
  \label{eq:lsq-ab}
\end{equation}
ここで $\sum$ は $i=1$ から $n$ までの和を表します。

% オームの法則（$V=RI$）のように，理論的に原点を通る比例関係になることがわかっている場合は，直線 $y=ax$ を当てはめます。このときの傾きは
% \begin{equation}
%   a = \frac{\sum x_i y_i}{\sum x_i^2}
%   \label{eq:lsq-origin}
% \end{equation}
% となります。

実験で扱う $x$ や $y$ は，電圧や電流などの単位をもつ物理量です。そのため，$y=ax+b$ を物理量の式として考えると，係数 $a$，$b$ 自体も単位（次元）をもちます。傾き $a$ は「$y$ の単位 ÷ $x$ の単位」，切片 $b$ は「$y$ の単位」をもつ量です。例えば，横軸が電流 $I$（A），縦軸が電圧 $V$（V）であれば，傾きの単位は V/A $=\Omega$ となり，傾きは抵抗を表します。\wave{傾きや切片を書くときは，数値だけでなく単位も必ず付けてください。}

\begin{hosoku}[Excel で求めるには]
\begin{itemize}
  \item 散布図に測定値をプロットし，「近似曲線の追加」で「線形近似」を選ぶと，最小二乗法による直線が引けます。「グラフに数式を表示する」にチェックを入れると，傾きと切片が表示されます。原点を通る直線にするときは，「切片」を0に設定します。
  \item 関数を使う場合は，傾きを \texttt{=SLOPE(}$y$ の範囲\texttt{,}$x$ の範囲\texttt{)}，切片を \texttt{=INTERCEPT(}$y$ の範囲\texttt{,}$x$ の範囲\texttt{)} で求めます。
  \item \wave{Excel では数値だけが表示されるので，単位を付けて読み取ります。}例えば電流を mA 単位で入力すると，傾きの単位は V/mA になります。$\SI{1}{V/mA}=\SI{1}{k\ohm}$ なので，表示された傾きが $-0.00074$ なら $a=-0.00074\ \mathrm{V/mA}=\SI{-0.74}{\ohm}$ です。
\end{itemize}
\end{hosoku}

% ------------------------------------------------------------
\section{計算例}

直流回路（テスター）の実験2では，乾電池から取り出す電流 $I$ と端子電圧 $V$ の関係が，$V=E-rI$（$E$：起電力，$r$：内部抵抗）の直線になると考えます。測定値を横軸 $I$，縦軸 $V$ として最小二乗法で直線 $V=aI+b$ を求めると，傾き $a$ から内部抵抗 $r=-a$ が，切片 $b$ から起電力 $E=b$ が求められます。

\begin{hosoku}[計算例]
ある乾電池について，抵抗値を\SI{1000}{\ohm}から\SI{100}{\ohm}まで変えて10組の $(I, V)$ を測定し，最小二乗法で直線 $V=aI+b$ を求めると，
\[
  a = \SI{-0.74}{\ohm},\qquad b = \SI{1.568}{V}
\]
となりました。傾き $a$ は電圧÷電流の次元をもつので単位は $\Omega$，切片 $b$ は電圧の次元をもつので単位は V です。したがって，内部抵抗は $r=-a=\SI{0.74}{\ohm}$，起電力は $E=b=\SI{1.568}{V}$ です。
\end{hosoku}

% % ------------------------------------------------------------
% \section{注意}

% \begin{itemize}
%   \item 最小二乗法を使う前に，必ず測定値をグラフにプロットし，測定点が直線に並んでいるかを確かめてください。\wave{直線にならない関係に無理に直線を当てはめても，意味のある結果は得られません。}
%   \item 他の点から大きく外れた測定点がある場合は，測定のやり直しや原因を検討してください。
%   \item 求めた傾きや切片も測定値から計算した値なので，有効数字を考えて書いてください。
% \end{itemize}
